\documentclass[11pt,a4paper]{article}
\usepackage{curso}
\begin{document}
\portada{Lección}{10.8}{Métricas de energía para modelos de lenguaje}

\noindent\textbf{Unidad 10.}\quad\textbf{Objetivos:} J/token, J/FLOP y MFU; titular frente a ganancia neta; el impuesto de la dispersión; cómo escribir una comparación justa.
\textbf{Prerrequisitos:} 10.1, 10.7.

\section{Métricas}
\[E_{\text{token}}=\frac{2P\,e_{\text{FLOP}}}{\text{MFU}}.\]
Con $P=11{,}55$M y $e_{\text{FLOP}}=0{,}177$ pJ (cifras del artículo de Z1T): $4{,}09\ \mu$J con MFU 100\,\%, $40{,}9\ \mu$J con 10\,\%: reconstruye exactamente la referencia
H100 de la tabla de Z1T.

\section{El titular \evP}
$139\times$ supone MFU 10\,\% (decode con lote 1). El propio artículo cita MFU $\approx40\,\%$ en modelos tipo Llama 3: con ese valor la razón es $\approx35\times$.

\section{Mismo alcance \evE}
El titular excluye la cabeza de vocabulario. Con $D=512$ y un vocabulario \emph{supuesto} de $V=50\,257$ (GPT-2; el artículo no lo da), la cabeza en la H100 cuesta
$2DV\,e_{\text{FLOP}}/\text{MFU}=9{,}1\ \mu$J$/\text{MFU}$.
\begin{center}\small\begin{tabular}{lrrr}\toprule
$\mu$J/token & cuerpo & cabeza & total\\\midrule
H100, MFU 10\,\% & 40{,}9 & 91 & 132\\ H100, MFU 40\,\% & 10{,}2 & 23 & 33\\ Z1T + FPGA & 0{,}29 & 136 & 137\\\bottomrule
\end{tabular}\end{center}
Razón con igual alcance: $\approx0{,}97$ (10\,\%) y $\approx0{,}24$ (40\,\%): empate o gana la GPU. La cabeza en la FPGA domina; moverla a un dispositivo eficiente es prioritario.

\section{Impuesto de la dispersión \evP}
A igual tamaño, el modelo disperso es peor; igualar a GPT-2 small exige $\sim10\times$ más FLOPs de entrenamiento. La penalización de tamaño en inferencia a igual calidad no se
reporta. Ganancia neta $=$ titular $/$ (factor de MFU $\cdot$ factor de alcance $\cdot$ impuesto de calidad).

\section{Comparación justa}
Declarar: (1) MFU y lote de la GPU; (2) alcance (cabeza, movimiento de datos, número de chips); (3) punto de calidad; (4) razones del sistema y solo del tejido; (5) etiqueta de
evidencia de cada número.

\begin{aplicacion}
Tabla titular frente a neto en el laboratorio L29.
\end{aplicacion}

\begin{error}
\begin{itemize}[nosep]
\item Comparar con distinto alcance. \quad Elegir el MFU que maximiza la razón.
\item Comparar modelos de distinta calidad. \quad Mezclar medidas y proyecciones sin etiquetas.
\end{itemize}
\end{error}

\begin{thebibliography}{9}\small
\bibitem{z1t} Extropic, Z1T, \url{https://extropic.ai/writing/z1t} (septiembre de 2026).
\bibitem{palm} A. Chowdhery et al., PaLM: scaling language modeling with Pathways, arXiv:2204.02311 (2022).
\bibitem{green} R. Schwartz et al., Green AI, \emph{Commun. ACM} 63, 54 (2020).
\end{thebibliography}
\end{document}
